% JAF 2026 contributed talk abstract (anonymous; ≤1 page excluding references)
% Submit as PDF with separate surname.txt (name, affiliation, email).
% https://jaf45warsaw.wfz.uw.edu.pl/

\documentclass[11pt]{article}
\usepackage[margin=1in]{geometry}
\usepackage{amsmath,amssymb}
\usepackage{microtype}

\begin{document}
\noindent
\textbf{Executable Self-Justifying Axiom Systems over U-Grounding Arithmetic:
Mechanizing Willard's Tableau Apparatus in a Weak-Arithmetic Setting}

\medskip
\noindent
Willard's self-justifying axiom systems (SJAS) form a family of weak
arithmetical theories that can consistently extend themselves with a
carefully formulated self-consistency principle, despite living in the
shadow of G\"odel's second incompleteness theorem.
The construction is not merely a model-theoretic curiosity: it depends on a
precise interaction between (i)~a \emph{U-grounding} arithmetic base, in
which addition is total but multiplication is available only as a graph
relation $\mathit{mult}(x,y,z)$ rather than a function symbol, and (ii)~a
reflected semantic-tableau proof apparatus~$D$, together with arithmetized
predicates for codes of formulae, proofs, and finite axiom bases~$\beta$.
In an $IS\#_D(\beta)$ system, the self-consistency sentence asserts that no
$D$-proof derives the code of a contradiction from the code of the finite
system $\beta$.
The boundary at which self-justification becomes possible is therefore
proof-theoretic and intensional: it is the failure of a particular encoded
tableau to close under the rules and size constraints Willard imposes, not
an extensional identification of provability with truth in a stronger
theory.

\smallskip
\noindent
This talk reports on an executable mechanization of the $IS\#_D(\beta)$ line
in \emph{Proflog}, a logic programming language whose operational semantics
is Fitting/Smullyan semantic tableaux with equality and subsidiary
procedure calls.
The implementation makes the weak-arithmetic object language, the reflected
clause basis, and the arithmetized proof predicate
$\mathsf{TabPrf}_\beta(s,t,p)$---``there is a semantic-tableau proof~$p$ of
theorem~$t$ in system~$s$''---first-class objects that can be queried,
executed, and audited at the proof-certificate level.
Because SJAS self-reference is sensitive to proof-tree branching, closure
rules, and encoded proof size (including Willard's conventional tableaux
encoding requirement), the mechanization is organized around \emph{visible}
tableau descent from source formulae to kernel proof terms, rather than
host-language evaluation of truth values.

\smallskip
\noindent
From the perspective of weak arithmetics, three aspects of the work may be
of particular interest.
First, the U-grounding design is treated as an executable signature and
coding discipline, not only as a metatheoretic hypothesis: numerals, bounded
formula classes ($\Delta^\star_0$, $\Pi^\star_1$, $\Sigma^\star_1$), and
Level-1 self-consistency schemata are generated and classified inside the
same first-order setting used for ordinary programs.
Second, the proof predicate is implemented as a structural checker over
encoded tableau trees, aligning with Willard's semantic-tableau proof
predicate for ordinary tableaux; ongoing correspondence work establishes a
biconditional over a sharply bounded first fragment of certificates, with
equality, procedure-call, and quantifier constructors shown unreachable in
that fragment.
Third, the executable setting supports \emph{computational} exploration of
Willard's ``strength dial'': the same tableau kernel runs Fitting-style
reference programs and SJAS-encoded proof obligations, showing where tableaux
close, remain open, or become unresolved when coding strength changes.

\smallskip
\noindent
The talk summarizes the weak-arithmetic commitments of the mechanized
$IS\#_D(\beta)$ builder, exhibits arithmetized proof checking on small
examples, and states what is not yet proved: full Willard equivalence,
U-grounding numeral decoding in the correspondence proof, and discharge of
$\beta$-validity.
The broader question---which weak-arithmetic and tableau features must be
preserved when a self-justifying system is made computational---is posed for
discussion.

\smallskip
\noindent
\textbf{Keywords:} weak arithmetic, U-grounding, self-justifying axiom
systems, semantic tableaux, proof coding, incompleteness.

\begin{thebibliography}{9}
\bibitem{willard2001}
D.~Willard,
Self-verifying axiom systems, the incompleteness theorem and related
reflection principles,
\emph{J.~Symbolic Logic} \textbf{66} (2001), 536--596.

\bibitem{willard2013}
D.~Willard,
The significance of $\Pi_1$-definable cutsets for incompleteness,
feasibility, and self-justification,
\emph{Ann. Pure Appl. Logic} \textbf{164} (2013), 902--912.

\bibitem{willard2004}
D.~Willard,
Addition, multiplication, and self-inspectable theories,
\emph{Ann. Pure Appl. Logic} \textbf{130} (2004), 27--41.

\bibitem{fitting1994}
M.~Fitting,
\emph{First-Order Logic and Automated Theorem Proving},
Springer, 1990/1996.
\end{thebibliography}

\end{document}
